\documentclass[11pt]{amsart}

\usepackage[T1]{fontenc}
\usepackage{lmodern}
\usepackage{microtype}
\usepackage{mathtools,amssymb,amsthm}
\usepackage[hidelinks]{hyperref}

\hypersetup{
  pdftitle={A primitive weird number with two square odd prime factors and Omega = 8}
}

\newtheorem{theorem}{Theorem}
\newtheorem{lemma}[theorem]{Lemma}
\newtheorem{corollary}[theorem]{Corollary}

\DeclareMathOperator{\Om}{\Omega}
\newcommand{\sig}{\sigma}

\emergencystretch=3em

\title[A primitive weird number with two square odd prime factors]
{A Primitive Weird Number with Two Square Odd Prime Factors
 and $\Omega=8$}

\date{}

\subjclass[2020]{11A25, 11B83, 11Y55}
\keywords{weird number, primitive weird number, primitive abundant number,
sum-of-divisors function, number of prime factors with multiplicity}

\begin{document}

\begin{abstract}
A \emph{weird} number is an abundant number that is not semiperfect, and a
\emph{primitive weird number} (PWN) is a weird number no proper divisor of
which is weird. For $n\in\mathbb N$ let $\Om_n$ denote the least value of
$\Om(m)$, the number of prime factors of $m$ counted with multiplicity, over
all PWN $m$ having exactly $n$ square odd prime factors. Open Question~1 of
Amato, Hasler, Melfi and Parton, quoted here from arXiv:1802.07178v2, records
$\Om_1=7$, $8\le\Om_2\le12$, $8\le\Om_3\le15$, and $\Om_n\ge8$ for $n\ge2$. We
exhibit
\[
 m=2^2\cdot13^2\cdot19^2\cdot46219\cdot1108619=12504224434300196 ,
\]
a PWN with exactly two square odd prime factors and $\Om(m)=8$, so
$\Om_2\le8$. With their lower bound $\Om_2\ge8$, which is quoted and not
reproved here, this gives $\Om_2=8$ and shows that their bound $\Om_n\ge8$ for
$n\ge2$ is attained; the equality is conditional on that quoted bound. The
verification is exact
integer arithmetic: two primality checks by trial division, then
$\sig(m)-2m=8$ with the divisors of $m$ that are at most $8$ being $1,2,4$,
whose total is $7$, and five comparisons showing the maximal divisors of $m$
deficient.
\end{abstract}

\maketitle

\section{The question}

Throughout, $\sig$ is the sum-of-divisors function, $\Om(m)$ is ``the number
of prime factors counted with their multiplicity'' and $\omega(m)$ is the
number of distinct prime factors. Write
\[
 \Delta(m)=\sig(m)-2m ,
\]
so that $m$ is \emph{abundant} exactly when $\Delta(m)>0$ and
\emph{deficient} exactly when $\Delta(m)<0$. Following Amato, Hasler, Melfi
and Parton \cite{AHMP19}, a number is \emph{semiperfect} if it is
expressible as a sum of distinct proper divisors of itself; a \emph{weird}
number is ``a number which is abundant but not semiperfect''; and a
\emph{primitive weird number} (PWN) is ``a weird number which is not a
multiple of any smaller weird number''. A \emph{square odd prime factor} of
$m$ is an odd prime $p$ with $p^2\mid m$.

In the section ``Open problems'' of \cite{AHMP19} the following appears, as
the first numbered Open Question of that section. It and the theorem below are
quoted from arXiv:1802.07178v2, which was not compared with the published
version.

\medskip
\noindent\textbf{Open Question 1} (\cite{AHMP19}).
\emph{For each $n$ in $\mathbb N$, find a PWN with exactly $n$ square odd
prime factors and the least $\Om=\Om_n$. From the previous section and
\textup{[the theorem quoted below]} we obtain $\Om_1=7$,
$8\le\Om_2\le12$, $8\le\Om_3\le15$, and in general if $n\ge2$ we have
$\Om_n\ge8$.}

\medskip
The lower bounds quoted there are the authors' own theorem, which reads
verbatim and in full: \emph{``There are no PWN $m$ with a quadratic
or higher power odd prime factor and $\Om(m)<7$. There are no PWN $m$ with $2$
quadratic odd prime factor and $\Om(m)=7$. There are no PWN $m$ with a cubic
or higher power odd prime factor and $\Om(m)=7$.''} The bound $\Om_2\ge8$ is
taken here from that theorem as a whole, transcribed and not reproved. We do
not apportion it among the three sentences: they distinguish a quadratic from a
cubic or higher odd prime factor, whereas the definition of square odd prime
factor used here, $p^2\mid m$, does not, and we do not settle which reading of
``quadratic odd prime factor'' the source intends. The upper bound
$\Om_2\le12$ is the smallest entry of the $\Om$ column of their table of
PWN with more than one square odd prime factor, which reads
$12,12,12,13,14,15$.

\section{The witness}

\begin{theorem}\label{thm:main}
Let
\[
 m=2^2\cdot13^2\cdot19^2\cdot46219\cdot1108619=12504224434300196 .
\]
Then $m$ is a primitive weird number, $\Om(m)=8$, and $m$ has exactly two
square odd prime factors, namely $13$ and $19$. Consequently $\Om_2\le8$, and
with the lower bound $\Om_2\ge8$ quoted from \cite{AHMP19} in \S1, $\Om_2=8$.
\end{theorem}

The following criterion is a proposition of \cite{AHMP19}, where it reads:
\emph{``An abundant number $n$ is weird iff $\Delta(n)$ cannot be expressed
as a sum of distinct proper divisors of $n$''}. A proof is included for
completeness.

\begin{lemma}\label{lem:crit}
An abundant number $N$ is weird if and only if $\Delta(N)$ is not a sum of
distinct proper divisors of $N$.
\end{lemma}

\begin{proof}
Let $P$ be the set of proper divisors of $N$, so
$\sum_{d\in P}d=\sig(N)-N=N+\Delta(N)$. If $S\subseteq P$ then its
complement $T=P\setminus S$ satisfies
$\sum_{d\in T}d=N+\Delta(N)-\sum_{d\in S}d$. Hence
$\sum_{d\in S}d=N$ holds for some $S\subseteq P$ if and only if
$\sum_{d\in T}d=\Delta(N)$ holds for some $T\subseteq P$. The left-hand
condition is semiperfection and the right-hand one is representability of
$\Delta(N)$; an abundant number is weird exactly when the former fails.
\end{proof}

\begin{proof}[Proof of Theorem~\ref{thm:main}]
\emph{(i) The factorisation.} $46219$ and $1108619$ are prime: neither is
divisible by any prime up to $\lfloor\sqrt{46219}\rfloor=214$, respectively
$\lfloor\sqrt{1108619}\rfloor=1052$. Hence
$\Om(m)=2+2+2+1+1=8$ and $\omega(m)=5$, and the odd primes occurring in $m$
to an exponent at least $2$ are exactly $13$ and $19$: two of them, both to
the exponent exactly $2$, so the two readings of ``square odd prime factor''
discussed in \S1 agree for $m$.

\emph{(ii) Abundance.} By multiplicativity,
\[
 \sig(m)=7\cdot183\cdot381\cdot46220\cdot1108620=25008448868600400 ,
\]
\[
 2m=25008448868600392 ,
\]
so
\[
 \Delta(m)=8>0 .
\]

\emph{(iii) Weird.} The smallest odd prime factor of $m$ is $13>8$, so the
divisors of $m$ that are at most $\Delta(m)=8$ are exactly $1$, $2$ and $4$,
and
\[
 1+2+4=7<8 .
\]
No sum of distinct divisors of $m$ can therefore equal $8$ --- the eligible
divisors do not even total $8$ --- so by Lemma~\ref{lem:crit} the abundant
number $m$ is weird.

\emph{(iv) Primitive abundant.} The abundancy $\sig(x)/x$ is non-decreasing
along divisibility, so every proper divisor of $m$ is deficient as soon as
the five maximal ones are, and these are settled by exact integer
comparison:
\[
\renewcommand{\arraystretch}{1.15}
\begin{array}{lrr}
 x & \sig(x)\hphantom{00} & 2x\hphantom{00}\\[2pt]
 \hline
 m/2       & 10717906657971600 & 12504224434300196\\
 m/13      & 1913214667543200  & 1923726836046184\\
 m/19      & 1312779468168000  & 1316234150978968\\
 m/46219   & 541074185820      & 541085892568\\
 m/1108619 & 22558179420       & 22558199768
\end{array}
\]
Each left entry is smaller than the right entry beside it, so each $m/q$ is
deficient and $m$ is primitive abundant. In particular no proper divisor of
$m$ is abundant, hence none is weird, and $m$ is a \emph{primitive} weird
number.

\emph{(v) The closure.} $m$ is a PWN with exactly two square odd prime
factors and $\Om(m)=8$, so $\Om_2\le8$; the theorem of \cite{AHMP19} quoted
in \S1 gives $\Om_2\ge8$; hence $\Om_2=8$.
\end{proof}

\begin{corollary}
Granting the bound $\Om_n\ge8$ recorded for $n\ge2$ in Open Question~$1$ of
\cite{AHMP19}, it is attained at $n=2$.
\end{corollary}

\medskip
\noindent\textbf{Status.}
Only the case $n=2$ of Open Question~1 is settled. The bound $\Om_2\le8$ is
proved above; the matching $\Om_2\ge8$ is the theorem of \cite{AHMP19} quoted
in \S1, transcribed and not reproved, so the equality $\Om_2=8$ is conditional
on that theorem. Nothing here bounds $\Om_n$ for $n\ge3$.

The integer itself is not new: $12504224434300196$ appears as the tenth term
of OEIS \texttt{A063788} \cite{OEIS}, ``numbers $k$ such that
$\sig(k)=2k+\Om(k)$'', consistently with $\Delta(m)=8=\Om(m)$ above. We make
no priority claim for the observation that $m$ is a PWN with two square odd
prime factors, nor for the closure that follows from it; we have not
established that either is unrecorded.

\begin{thebibliography}{9}

\bibitem{AHMP19}
G.~Amato, M.~F. Hasler, G.~Melfi, and M.~Parton,
\emph{Primitive abundant and weird numbers with many prime factors},
J. Number Theory \textbf{201} (2019), 436--459.
\href{https://doi.org/10.1016/j.jnt.2019.02.027}{doi:10.1016/j.jnt.2019.02.027};
\href{https://arxiv.org/abs/1802.07178v2}{arXiv:1802.07178v2}.

\bibitem{OEIS}
OEIS Foundation Inc., \emph{The On-Line Encyclopedia of Integer Sequences}:
sequence \texttt{A063788}, numbers $k$ with $\sigma(k)=2k+\Omega(k)$,
\href{https://oeis.org/A063788}{oeis.org/A063788}.

\end{thebibliography}

\end{document}
